\documentclass[UTF8,zihao=5,twocolumn]{ctexart}

\usepackage[a4paper,top=1.8cm,bottom=1.9cm,left=1.55cm,right=1.55cm]{geometry}
\usepackage{booktabs}
\usepackage{caption}
\usepackage{hyperref}
\usepackage{listings}
\usepackage{xcolor}
\usepackage{array}

\setlength{\columnsep}{0.65cm}
\setlength{\parindent}{2em}
\setlength{\parskip}{0pt}
\linespread{1.05}
\pagestyle{plain}
\captionsetup{font=small,labelfont=bf}

\ctexset{
  section = {
    format = \large\bfseries,
    name = {},
    number = \arabic{section},
    aftername = \quad
  },
  subsection = {
    format = \normalsize\bfseries,
    name = {},
    number = \arabic{section}.\arabic{subsection},
    aftername = \quad
  }
}

\hypersetup{
  colorlinks=true,
  linkcolor=black,
  citecolor=black,
  urlcolor=blue,
  pdftitle={基于 TCP Socket 与 select 的多人聊天室系统},
  pdfauthor={Linux 高级编程课程设计}
}

\lstdefinestyle{cppstyle}{
  language=C++,
  basicstyle=\ttfamily\scriptsize,
  keywordstyle=\bfseries\color{blue!55!black},
  commentstyle=\itshape\color{green!35!black},
  stringstyle=\color{red!55!black},
  numbers=left,
  numberstyle=\tiny\color{gray},
  stepnumber=1,
  numbersep=5pt,
  breaklines=true,
  columns=fullflexible,
  keepspaces=true,
  frame=single,
  framerule=0.25pt,
  rulecolor=\color{black!25},
  tabsize=2,
  showstringspaces=false,
  captionpos=b
}

\lstdefinestyle{textstyle}{
  basicstyle=\ttfamily\scriptsize,
  breaklines=true,
  columns=fullflexible,
  keepspaces=true,
  frame=single,
  framerule=0.25pt,
  rulecolor=\color{black!25},
  showstringspaces=false
}

\newcommand{\code}[1]{\texttt{#1}}
\newcommand{\sysname}{TCP Select Chatroom}

\begin{document}

\twocolumn[
\begin{center}
{\LARGE \bfseries 基于 TCP Socket 与 select 的多人聊天室系统}\\
\vspace{0.7em}
{\normalsize 《Linux 高级编程》课程设计报告}\\
\vspace{0.5em}
{\small 作者：\underline{\hspace{3.5cm}} \quad 学号：\underline{\hspace{3.5cm}} \quad 日期：2026 年 7 月}\\
\vspace{1.1em}
\end{center}

\begin{abstract}
\noindent
本文设计并实现了一个运行在 Linux 系统上的多人聊天室系统。系统使用 C++ 编写，基于 TCP Socket 完成客户端与服务端之间的可靠通信，服务端采用单进程 \code{select} I/O 多路复用模型同时管理多个客户端连接。客户端连接服务端后使用昵称登录，支持群聊广播、用户私聊、在线用户列表查询、正常退出、异常断开处理和服务端日志记录。项目重点体现 Linux 网络编程中的文件描述符管理、Socket 生命周期、文本协议解析、TCP 字节流拆包、异常清理和可验证测试。实验结果表明，该系统能够在 Linux 环境下完成多客户端并发聊天，并能在非法命令、重复昵称、超长消息和连接异常等情况下保持服务端可用。
\end{abstract}

{\small\noindent\textbf{关键词：} Linux；TCP Socket；select；I/O 多路复用；多人聊天室；C++}\\
\vspace{1.2em}
]

\section{引言}

《Linux 高级编程》课程设计要求在 Linux 系统上实现一个具有服务端和客户端的软件系统，并体现 Socket 编程或进程通信等高级编程内容。多人聊天室是典型的网络类应用：服务端需要监听客户端连接，接收某个客户端发送的消息，并将消息分发给其他客户端；客户端需要连接服务端，发送用户输入，并实时接收服务端转发的消息。

本文实现的 \sysname{} 选择 TCP 作为传输协议，使用 \code{select} 作为 I/O 多路复用机制。相比为每个客户端创建线程，单进程 \code{select} 模型更适合展示 Linux 文件描述符和事件循环的核心思想，也便于在课程设计规模内控制代码复杂度。

本项目的目标包括：
\begin{itemize}
  \item 实现一个 TCP 聊天室服务端和一个 TCP 聊天室客户端；
  \item 服务端使用单进程加 \code{select} 管理多个客户端连接；
  \item 支持昵称登录、群聊广播、私聊、在线用户列表和退出；
  \item 处理重复昵称、非法命令、超长消息和异常断开；
  \item 将启动、登录、退出、异常断开、群聊、私聊和错误事件写入日志。
\end{itemize}

为保持范围清晰，本系统不实现账号密码认证、数据库存储、离线消息、历史记录查询、图形界面和管理员踢人功能。

\section{系统设计}

\subsection{总体架构}

系统由两个可执行程序组成：\code{server} 和 \code{client}。服务端负责监听端口、接受连接、维护在线用户、解析命令、分发消息和写入日志；客户端负责连接服务端、发送登录命令、读取用户输入、接收服务端消息并在终端显示。

图~\ref{fig:architecture} 给出系统结构。协议层提供命令解析、响应生成和按行拆包；服务端状态层维护 fd、昵称和登录状态；服务端业务层负责连接管理与消息路由；客户端交互层将用户输入转换成协议命令，并将服务端响应格式化为可读文本。

\begin{figure*}[t]
\centering
\fbox{
\begin{minipage}{0.92\textwidth}
\centering
\small
\begin{tabular}{>{\centering\arraybackslash}p{0.18\textwidth}
                >{\centering\arraybackslash}p{0.20\textwidth}
                >{\centering\arraybackslash}p{0.20\textwidth}
                >{\centering\arraybackslash}p{0.18\textwidth}}
客户端终端 & \code{ChatClient} & TCP 文本协议 & \code{ChatServer} \\
\midrule
用户输入 & \code{InputHandler} & \begin{tabular}{@{}c@{}}\code{LOGIN}/\code{MSG}/\code{PRIVATE}\\\code{LIST}/\code{QUIT}\end{tabular} & \code{ClientManager} \\
消息显示 & \code{Display} & \begin{tabular}{@{}c@{}}\code{OK}/\code{ERROR}\\\code{CHAT}/\code{SYSTEM}\end{tabular} & \code{Logger}
\end{tabular}
\end{minipage}
}
\caption{系统总体架构}
\label{fig:architecture}
\end{figure*}

\subsection{代码组织}

项目代码位于 \code{code/} 目录，主要结构如表~\ref{tab:files} 所示。

\begin{table}[h]
\centering
\caption{主要文件与职责}
\label{tab:files}
\small
\begin{tabular}{p{0.45\linewidth}p{0.43\linewidth}}
\toprule
文件 & 作用 \\
\midrule
\code{protocol.*} & 协议解析、响应生成、按行缓冲 \\
\code{server.*} & 服务端监听、\code{select} 循环、消息路由 \\
\code{client\_manager.*} & 在线用户状态管理 \\
\code{logger.*} & 服务端运行日志记录 \\
\code{client.*} & 客户端连接、登录和交互循环 \\
\code{input\_handler.*} & 用户输入到协议命令的转换 \\
\code{display.*} & 服务端响应的终端显示格式化 \\
\bottomrule
\end{tabular}
\end{table}

\section{TCP Socket 通信}

TCP 是面向连接、可靠、有序的字节流协议。聊天室要求客户端和服务端之间保持长连接，并保证消息按顺序到达，因此使用 TCP 作为传输层协议。

服务端启动时依次调用 \code{socket}、\code{setsockopt}、\code{bind} 和 \code{listen}。当监听 fd 可读时，调用 \code{accept} 接收新客户端。核心代码如代码~\ref{lst:server-start} 所示。

\begin{lstlisting}[style=cppstyle,caption={服务端 Socket 初始化},label={lst:server-start}]
listen_fd_ = socket(AF_INET, SOCK_STREAM, 0);
setsockopt(listen_fd_, SOL_SOCKET, SO_REUSEADDR,
           &reuse, sizeof(reuse));
bind(listen_fd_, reinterpret_cast<sockaddr*>(&address),
     sizeof(address));
listen(listen_fd_, kListenBacklog);
\end{lstlisting}

客户端启动后解析服务端 IP、端口和昵称，创建 TCP Socket，再调用 \code{connect} 建立连接。连接成功后立即发送 \code{LOGIN <nickname>}，服务端返回 \code{OK} 后客户端进入交互循环。

由于 TCP 是字节流协议，一次 \code{recv} 可能读到半条消息，也可能读到多条消息。因此系统约定每条协议消息以换行符结束，并为每个连接维护 \code{LineBuffer}。该缓冲区将新收到的字节追加到内部字符串，遇到换行符时输出完整消息，未形成完整行的部分继续保留。

\section{select I/O 多路复用}

\subsection{服务端事件循环}

服务端采用单进程 \code{select} 模型。每次进入事件循环时，服务端重新构造 \code{fd\_set}，将监听 fd 和所有客户端 fd 加入集合，然后调用 \code{select} 等待可读事件。若监听 fd 可读，说明有新连接；若客户端 fd 可读，说明该客户端发送了消息或已经断开。

\begin{lstlisting}[style=cppstyle,caption={服务端 select 主循环片段},label={lst:server-select}]
FD_ZERO(&read_fds);
FD_SET(listen_fd_, &read_fds);
int max_fd = listen_fd_;

for (int fd : clients_.fds()) {
    FD_SET(fd, &read_fds);
    if (fd > max_fd) {
        max_fd = fd;
    }
}

int ready = select(max_fd + 1, &read_fds,
                   nullptr, nullptr, timeout_ptr);
\end{lstlisting}

该模型的优势在于：服务端不需要为每个客户端创建线程，所有连接通过 fd 集合统一调度，资源管理逻辑集中，适合课程设计中展示 Linux 文件描述符和 I/O 多路复用机制。

\subsection{客户端事件循环}

客户端也使用 \code{select} 同时监听标准输入和服务端 Socket。标准输入可读时，客户端读取用户命令并发送协议消息；服务端 Socket 可读时，客户端读取服务端响应并格式化输出。这样可以在不使用线程的情况下实现“边输入边接收”。

\section{通信协议}

系统使用按行分隔的文本协议。文本协议便于调试，也便于使用终端或脚本模拟客户端。客户端命令和服务端响应如表~\ref{tab:protocol} 所示。

\begin{table}[h]
\centering
\caption{聊天室文本协议}
\label{tab:protocol}
\small
\begin{tabular}{ll}
\toprule
方向 & 消息格式 \\
\midrule
客户端 & \code{LOGIN <nickname>} \\
客户端 & \code{MSG <content>} \\
客户端 & \code{PRIVATE <target> <content>} \\
客户端 & \code{LIST} \\
客户端 & \code{QUIT} \\
\midrule
服务端 & \code{OK <message>} \\
服务端 & \code{ERROR <reason>} \\
服务端 & \code{SYSTEM <message>} \\
服务端 & \code{CHAT <nickname> <content>} \\
服务端 & \code{PRIVATE <from> <content>} \\
服务端 & \code{USERLIST <name1>,<name2>,...} \\
\bottomrule
\end{tabular}
\end{table}

协议约束包括：客户端连接后必须先发送 \code{LOGIN}；昵称长度为 1 到 20 个字符，只允许英文字母、数字和下划线；单条协议消息最大 1024 字节；未登录客户端只能发送 \code{LOGIN} 或 \code{QUIT}。

\begin{lstlisting}[style=cppstyle,caption={昵称合法性检查},label={lst:nickname}]
bool is_valid_nickname(const std::string& nickname) {
    if (nickname.empty() ||
        nickname.size() > kMaxNicknameLength) {
        return false;
    }

    for (char c : nickname) {
        if (c != '_' && !is_ascii_letter_or_digit(c)) {
            return false;
        }
    }
    return true;
}
\end{lstlisting}

\section{服务端实现}

\subsection{客户端状态管理}

服务端使用 \code{ClientManager} 管理客户端状态。每个客户端记录 fd、昵称、登录状态和接收缓冲区。状态管理同时支持按 fd 查找和按昵称查找：按 fd 查找用于处理 Socket 事件，按昵称查找用于私聊和重复昵称检测。

客户端连接后先处于未登录状态。服务端只允许未登录客户端发送 \code{LOGIN} 或 \code{QUIT}。收到 \code{LOGIN} 后，服务端检查昵称是否合法且未被占用，再写入昵称和登录状态。登录成功后，服务端向该客户端返回 \code{OK login successful}，并向其他在线用户广播上线通知。

\subsection{消息路由}

群聊消息由 \code{handle\_group\_message} 处理。服务端确认发送者已登录后，生成 \code{CHAT <nickname> <content>} 响应，并遍历所有已登录客户端，将消息发送给除发送者以外的用户。

私聊消息由 \code{handle\_private\_message} 处理。服务端根据目标昵称查找目标客户端；目标不存在时返回 \code{ERROR user not online}；目标存在时仅向目标 fd 发送 \code{PRIVATE <from> <content>}，不会泄露给其他客户端。

\begin{lstlisting}[style=cppstyle,caption={私聊消息路由片段},label={lst:private}]
const ClientInfo* target_client =
    clients_.find_by_nickname(target);
if (target_client == nullptr ||
    !target_client->logged_in) {
    send_response(fd, make_error("user not online"));
    return;
}

send_response(target_client->fd,
              make_private(sender->nickname, content));
\end{lstlisting}

\subsection{异常处理与日志}

服务端需要处理重复昵称、非法昵称、未登录命令、非法命令、私聊目标不存在、超长消息、客户端异常断开、\code{send} 或 \code{recv} 失败等情况。对每个异常，服务端尽量只清理受影响连接，不影响其他在线用户。

\code{Logger} 将关键事件追加写入 \code{server.log}，事件类型包括 \code{STARTUP}、\code{LOGIN}、\code{QUIT}、\code{ABNORMAL\_DISCONNECT}、\code{GROUP\_CHAT}、\code{PRIVATE\_CHAT} 和 \code{ERROR}。日志写入失败时，服务端只向标准错误输出提示，不因为日志问题终止。

\section{客户端实现}

\subsection{启动与登录}

客户端启动格式如下：

\begin{lstlisting}[style=textstyle,caption={客户端启动命令}]
./client 127.0.0.1 8888 alice
\end{lstlisting}

客户端解析服务端 IP、端口和昵称，连接成功后发送 \code{LOGIN alice}。如果服务端返回 \code{OK}，客户端进入交互循环；如果返回 \code{ERROR} 或关闭连接，客户端输出失败原因并退出。

\subsection{用户输入处理}

\code{InputHandler} 将用户输入转换为协议命令，如表~\ref{tab:client-command} 所示。空输入会被忽略；\code{/w} 缺少目标或消息内容时，客户端在本地提示错误，不向服务端发送非法命令。

\begin{table}[h]
\centering
\caption{客户端输入映射}
\label{tab:client-command}
\small
\begin{tabular}{ll}
\toprule
用户输入 & 协议命令 \\
\midrule
\code{hello} & \code{MSG hello} \\
\code{/w bob hello} & \code{PRIVATE bob hello} \\
\code{/list} & \code{LIST} \\
\code{/quit} & \code{QUIT} \\
\bottomrule
\end{tabular}
\end{table}

\subsection{消息显示}

服务端响应是协议文本，直接显示可读性较差。客户端使用 \code{Display} 将响应转换为终端友好的格式。例如，\code{CHAT alice hello} 显示为 \code{[group] alice: hello}，\code{PRIVATE bob hi} 显示为 \code{[private] bob: hi}，\code{ERROR user not online} 显示为 \code{Error: user not online}。

\section{实验与验证}

\subsection{构建方法}

项目提供 \code{Makefile}。在 \code{code/} 目录执行 \code{make} 后生成 \code{server} 和 \code{client} 两个可执行文件。执行 \code{make clean} 可清理构建产物。

\begin{lstlisting}[style=textstyle,caption={构建与运行命令}]
cd code
make
./server 8888
./client 127.0.0.1 8888 alice
\end{lstlisting}

\subsection{测试覆盖}

本项目包含自动化测试和端到端手动测试。自动化测试覆盖协议解析、客户端状态管理、日志写入、客户端输入转换、显示格式化、客户端连接、客户端 \code{select} 循环、服务端监听、登录生命周期、消息路由和错误处理。主要测试结果如表~\ref{tab:tests} 所示。

\begin{table}[h]
\centering
\caption{测试覆盖摘要}
\label{tab:tests}
\small
\begin{tabular}{p{0.43\linewidth}p{0.47\linewidth}}
\toprule
测试文件 & 覆盖内容 \\
\midrule
\code{protocol\_tests} & 拆包、昵称、命令、1024 字节限制 \\
\code{server\_routing\_tests} & 群聊、私聊、在线列表、日志 \\
\code{server\_lifecycle\_tests} & 登录、退出、异常断开 \\
\code{server\_hardening\_tests} & 非法命令、超长消息、send/recv 失败 \\
\begin{tabular}{@{}l@{}}\code{client\_select}\\\code{\_loop\_tests}\end{tabular} & 标准输入和服务端 Socket 并发处理 \\
\bottomrule
\end{tabular}
\end{table}

重新整理本文时，已执行 \code{make -C code clean \&\& make -C code}，服务端和客户端均能成功编译。随后编译并运行 11 个 C++ 测试二进制，测试进程均以退出码 0 结束。端到端验证还检查了三客户端登录、群聊广播、私聊只到目标、在线列表、重复昵称拒绝、非法命令、超长消息、正常退出、异常断开和日志记录。

\section{结论}

本文实现了一个基于 Linux TCP Socket 和 \code{select} 的多人聊天室系统。系统通过按行文本协议实现客户端和服务端协作，通过 \code{ClientManager} 维护在线用户状态，通过服务端消息路由逻辑完成群聊、私聊和在线列表，通过 \code{Logger} 记录关键运行事件，并通过客户端 \code{select} 循环实现用户输入和服务端消息的并发处理。

课程设计表明，\code{select} 能够以较低复杂度支持多个客户端连接，并能清晰展示 Linux 文件描述符、Socket 生命周期和 I/O 多路复用机制。后续可以扩展账号认证、持久化聊天记录、图形界面、管理员功能，或者将服务端事件模型替换为更适合大量连接的 \code{epoll}。

\begin{thebibliography}{9}

\bibitem{unp}
W. Richard Stevens, Bill Fenner, and Andrew M. Rudoff.
\newblock \emph{UNIX Network Programming, Volume 1: The Sockets Networking API}.
\newblock Addison-Wesley, 3rd edition, 2003.

\bibitem{tlpi}
Michael Kerrisk.
\newblock \emph{The Linux Programming Interface}.
\newblock No Starch Press, 2010.

\bibitem{posix}
IEEE and The Open Group.
\newblock \emph{The Open Group Base Specifications Issue 7, POSIX.1-2017}.
\newblock 2018 edition.

\bibitem{manpages}
Linux man-pages project.
\newblock Manual pages for \code{socket(2)}, \code{select(2)}, \code{recv(2)} and \code{send(2)}.

\end{thebibliography}

\end{document}
